
\author[Mike Murphy]{Michael J. Murphy}
\institute[UiT]{UiT The Arctic Universtity of Norway}

% Font encoding
\usepackage[utf8]{inputenc}
\usepackage[T1]{fontenc}

% Remove non-deterministic metadata, for reproducible PDF builds
\usepackage[nopdftrailerid=true]{pdfprivacy}

% Use UiT's link colors in article mode too.
% (In presentation mode, this is handled by the UiT color theme.)
\mode<article>{
    \uitlinkcolors
}

\begin{document}

% Special title page setup for article mode
\mode<article>{
    % Render normal article title.
    \maketitle

    % Tell Beamer where to find the slides PDF to extract slides from.
    \setjobnamebeamerversion{uitbeamertest}

    % Draw the title slide as an image.
    \begin{center}
        \includeslide[width=\textwidth]{title-pic}
    \end{center}

    % Render table of contents.
    \tableofcontents
}

% Basic title frame
\begin{frame}[plain]
    \titlepage
\end{frame}

% Title frame with graphic
\begin{frame}[plain]
    \titlegraphic{\includegraphics
            [height=\paperheight, trim=0 0 1cm 0]
            {download/csiro-five-antennas-bg.jpg}}
            % (The trim key's arguments are left bottom right top)
    \tikzset{split pos = .50}
    \titlepage
    \label{title-pic}
\end{frame}

\begin{frame}{Antenna Photo Information}
    \centering
    \includegraphics[height=.25\paperheight]{download/csiro-five-antennas-bg.jpg}

    The antenna image used throughout this presentation is by
    John Masterson,
    \href{https://www.scienceimage.csiro.au/pages/about/}{CSIRO},
    licensed under \href{https://creativecommons.org/licenses/by/3.0/}{CC-BY-3.0}.
    It is used unmodified except for cropping or adding a color overlay.
    More info about the photo at \href{https://commons.wikimedia.org/wiki/File:CSIRO_ScienceImage_3881_Five_Antennas_at_Narrabri.jpg}{Wikimedia Commons}.
\end{frame}

\section{Color Tests}

\subsection{Basic Colors}

\begin{frame}{Color Test: Base UiT Palette}
    \scriptsize
    \uitcolortest

    \medskip

    \uitcolortest[rgb,hsb,gray]
\end{frame}

\subsection{Beamer Colors}
\label{beamercolortest}

\begin{frame}{Color Test: Beamer Colors}
    \begin{columns}
        \begin{column}{0.65\textwidth}
            \scriptsize
            \uitcolortestbeamerpalette
            \smallskip
        \end{column}
        \begin{column}{0.25\textwidth}
            \begin{block}{Test Block}
                Test block
            \end{block}
            \begin{alertblock}{Alert Block}
                Alert block
            \end{alertblock}
            \begin{exampleblock}{Example Block}
                Example block
            \end{exampleblock}
        \end{column}
    \end{columns}
\end{frame}

\begin{frame}{Color Test: Hyperlinks}

    \begingroup
        \scriptsize
        \begin{testcolors}[cmyk,RGB,HTML]
            \testcolor{uitlinkcolor}
            \testcolor{uiturlcolor}
        \end{testcolors}
    \endgroup

    \medskip

    \begin{itemize}
        \item Example internal ref to \hyperref[beamercolortest]{Beamer Color Test}
        \item Example URL \url{https://www.ctan.org/pkg/beamer}
        \item Example citation \cite{aa}
        \item Example file ref \url{uit.sty}
    \end{itemize}

    \medskip

    \begin{block}{The Bibliography Block}
        \begin{thebibliography}{1}
            \bibitem{aa} Jane Doe. "An example bibliography entry."
                \url{https://www.ctan.org/pkg/beamer}
        \end{thebibliography}
    \end{block}
\end{frame}

\begin{frame}{Angle Test}
    \begin{itemize}
        \def\ew{3cm}
        \def\eh{2.5cm}
        \item The split can be drawn with different parameters

            \tikz{\uitsplitbackground{0,0}{\ew, \eh}}
            \tikz{\uitsplitbackground[
                    split pos = .25,
                    split right/.append style = {color=UitRed, opacity=.4},
                ]{0,0}{\ew, \eh}}
            \tikz{\uitsplitbackground[
                    opacity=.6,
                    line width = 1em,
                ]{0,0}{\ew, \eh}}
            \tikz{\uitsplitbackground[split rays=true]{0,0}{\ew, \eh}}

        \item The split can be placed with transparency, as seen in
            \href{https://uit.no/ansatte/kommunikasjonstips/grafisk/maler/sub?p_document_id=861319}{UiT's example report covers}

            \begin{tikzpicture}
                \node[inner sep=0, outer sep=0] (bg)
                    {\includegraphics[height=\eh]{download/csiro-five-antennas-bg.jpg}};
                \uitsplitbackground[
                    split pos = .55,
                    split left/.style = {},
                    split right/.append style = {color=UitRed, opacity=.4},
                    ]{bg.south west}{bg.north east}
            \end{tikzpicture}
            \begin{tikzpicture}
                \node[inner sep=0, outer sep=0] (bg)
                    {\includegraphics[height=\eh]{download/csiro-five-antennas-bg.jpg}};
                \uitsplitbackground[
                    split pos = .75,
                    split left/.style = {},
                    split right/.append style = {color=UitYellow, opacity=.4},
                    ]{bg.south west}{bg.north east}
            \end{tikzpicture}
            \begin{tikzpicture}
                \node[inner sep=0, outer sep=0] (bg)
                    {\includegraphics[height=\eh]{download/csiro-five-antennas-bg.jpg}};
                \uitsplitbackground[
                    split pos=.33,
                    split rays=true,
                    split right/.style = {},
                    split ray right/.style = {},
                    ]{bg.south west}{bg.north east}
            \end{tikzpicture}
    \end{itemize}
\end{frame}

\begin{frame}{Rays Test}
    \begin{itemize}
        \item Here are the rays by themselves, without being cut off by the background.\\
            The scale on the 4th shape shows the origin and height parameters,\\
            which are based on the title page.
            \medskip

            \tikz{\uitsplitrays{0,0}{4em}}
            \tikz{\uitsplitrays[split ray/.append style={color=UitYellow}]{0,0}{4em}}
            \tikz{\uitsplitrays[split ray/.append style={color=UitRed}]{0,0}{4em}}
            \begin{tikzpicture}
                \uitsplitrays{0,0}{4em}
                \path[draw] (-2pt,0) -- (2pt,0)
                            (0,0) -- (0,4em)
                            (-2pt,4em) -- (2pt,4em);
                \path[fill] (0,0) circle (1pt);
            \end{tikzpicture}
    \end{itemize}
\end{frame}

\begin{frame}{Regular frame with background rays}
    To add decorative rays in background, \\
    use the \texttt{bgelement} option when including the \texttt{uit}
    inner theme.
\end{frame}

\end{document}
